\begin{definition}
    \textbf{Serie divergente a $+\infty$}

    Si dice che una successione $\left\{x_n\right\}_{n\in\mathbb{N}}$ diverge a $+\infty$
    se e solo se
    \[
        \forall M > 0 \exists n_0 \mid \forall n \geq n_0, x_n > M.
    \]
\end{definition}

\begin{definition}
    \textbf{Serie divergente a $-\infty$}

    Si dice che una successione $\left\{x_n\right\}_{n\in\mathbb{N}}$ diverge a $-\infty$
    se e solo se
    \[
        \forall m < 0 \exists n_0 \mid \forall n \geq n_0, x_n < m.
    \]
\end{definition}

\begin{definition}
    \textbf{Limite infinito di una successione}

    Si dice che il limite di una successione (il valore a cui tende per $n\rightarrow
    +\infty$) è infinito se e solo se la successione è divergente a $+\infty$ o $-\infty$.

    Se il limite di una successione è infinito, avremo che il limite di quella successione
    coinciderà con il valore di divergenza e scriveremo rispettivamente per una serie
    divergente a $+\infty$ e per una divergente a $-\infty$ che
    \[
        \underset{n\rightarrow+\infty}{\lim} x_n = +\infty
    \]
    e che
    \[
        \underset{n\rightarrow+\infty}{\lim} x_n = -\infty.
    \]

\end{definition}

\begin{example}\leavevmode

    \begin{enumerate}[label=(\roman*)]
        \item \[
            n \rightarrow +\infty \left(\underset{n\rightarrow+\infty}{\lim}n=+\infty\right)
        \]
        \item \[
            n^2 \rightarrow +\infty\left(\underset{n\rightarrow+\infty}{\lim}n^2=+\infty\right)
        \]
        \item \[
            -n^2 \rightarrow -\infty\left(\underset{n\rightarrow+\infty}{\lim}-n^2=-\infty\right)
        \]
        \item \[
            \frac{n^2}{n+1} \rightarrow +\infty
            \left(\underset{n\rightarrow+\infty}{\lim}\left[\frac{n^2}{n+1}\right]=+\infty\right)
        \]
    \end{enumerate}
\end{example}